\chapter{The Forbidden Subgraph, Resolved at \texorpdfstring{$d=2$}{d=2}}
\label{ch:forbidden-subgraph}

Chapter~\ref{ch:minimal-witnesses}'s third open question asked for a finite
family of forbidden context subgraphs characterizing real-realizability.
This chapter answers it completely at $d=2$: the forbidden family has
exactly one member, a bare cycle, and the characterization is that $G$ is
always real-realizable if and only if $G$ is a forest.

\section{Subdivision Preserves the Obstruction}

\begin{lemma}[Subdivision lemma]
\label{lem:subdivision}
\index{Subdivision lemma (lemma)}
Let $G'$ be obtained from a context graph $G$ by subdividing an edge
$(u,v)$ — replacing it with a path $u,w_1,\dots,w_k,v$ through new
vertices — and labeling every new edge with canonical angle $0$. Then $G$
is realizable with $(u,v)$ labeled $\theta$, together with the rest of
$G$'s labeling, if and only if $G'$ is realizable with the same labeling
on the unchanged edges and every subdivision edge labeled $0$.
\end{lemma}

\begin{proof}
A canonical angle of $0$ on an edge $(w_{i-1},w_i)$ is achieved exactly by
the real frame choice $a_{w_i} = a_{w_{i-1}}$ (one valid branch of the
fibre described in \S\ref{sec:gauge-single-frame}). Propagating this
choice along the entire subdivision path forces $a_v = a_u$ exactly,
identically to what a direct edge $(u,v)$ with canonical angle $0$ would
force. For any \emph{other} target angle $\theta$ on the original edge
$(u,v)$, choose instead $a_{w_1}=\cdots=a_{w_k}=a_u$ (still trivial edges)
and require $a_v - a_u$ to realize $\theta$ exactly as in the unsubdivided
graph — the trivial path contributes nothing to the relationship between
$a_u$ and $a_v$, so realizability of $G$ with $(u,v)=\theta$ and
realizability of $G'$ with the subdivision (trivial edges) plus $(u,v)$'s
role now played by the path's endpoints are the same condition on the same
free parameters $\{a_x\}_{x\in V(G)}$, with the new vertices'
angles pinned to $a_u$ and contributing no additional freedom or
constraint.
\end{proof}

\begin{corollary}[Every cycle length inherits the triangle obstruction]
\label{cor:cycle-inherits}
\index{Every cycle length inherits the triangle obstruction (corollary)}
For every $n\ge3$, there is an assignment of canonical angles to the edges
of the $n$-cycle $C_n$, individually admissible on every edge, that is not
globally real-realizable.
\end{corollary}

\begin{proof}
$C_n$ is a subdivision of $C_3$ for every $n\ge3$ (subdivide $n-3$ of the
triangle's edges, distributed arbitrarily, to reach length $n$). Apply
Lemma~\ref{lem:subdivision} to the mutually-unbiased triangle labeling of
Corollary~\ref{cor:real-mub} (canonical angle $\pi/4$ on all three original
edges), labeling every subdivision edge $0$: by
Corollary~\ref{cor:real-mub} the triangle labeling is unrealizable, so by
the lemma the resulting $C_n$ labeling is unrealizable too.
\end{proof}

This was checked computationally before being stated as a corollary: exact
branch-search confirmed the induced labeling is realizable on
neither $C_4$ nor $C_5$ (zero solutions found over all $4^3=64$ and
$4^4=256$ branch combinations respectively), while a control labeling
differing only in the closing edge (replacing the inherited $\pi/4$ with
the trivially-realizable $0$) is realizable on both.

\section{The Forbidden Subgraph Theorem}

\begin{theorem}[Forbidden subgraph, $d=2$]
\label{thm:forbidden-subgraph}
\index{Forbidden subgraph, d=2 (theorem)}
$\mathcal R_\R(G,2) = \mathcal D(G,2)$ (every individually admissible
edge-labeling of $G$ is globally real-realizable) if and only if $G$ is a
forest.
\end{theorem}

\begin{proof}
($\Leftarrow$) If $G$ is a forest, apply
Proposition~\ref{prop:tree-no-constraint} to each connected component: any
individually admissible labeling extends to a global real frame assignment
on each tree, hence on all of $G$.

($\Rightarrow$) If $G$ is not a forest, $G$ contains a cycle (a graph with
$|E|\ge|V|$ on some connected component must contain one, since a forest
has at most $|V|-1$ edges per component); let $C_n$, $n\ge3$, be that
cycle, viewed as a subgraph of $G$. Apply
Corollary~\ref{cor:cycle-inherits} to label $C_n$'s edges with an
individually-admissible, globally-unrealizable pattern, and label every
edge of $G$ outside $C_n$ with canonical angle $0$ (trivially admissible).
The resulting labeling of all of $G$ is individually admissible on every
edge; it cannot be globally realized, since any global realization would
in particular realize $C_n$'s edges, contradicting
Corollary~\ref{cor:cycle-inherits}. Hence $\mathcal D(G,2) \not\subseteq
\mathcal R_\R(G,2)$.
\end{proof}

\begin{remark}
This is the forbidden-subgraph family in its simplest possible form: a
single forbidden pattern, "any cycle," rather than a genuinely infinite or
complicated family. Equivalently: $\mathcal F = \{C_3, C_4, C_5,\dots\}$
works, but by Corollary~\ref{cor:cycle-inherits} it can be stated purely
topologically as $\mathcal F = \{\text{any cycle}\}$, i.e.\ $G$ is
forbidden-subgraph-free exactly when $\beta_1(G)=0$.
\end{remark}

\section{Contrast with \texorpdfstring{$d\ge3$}{d>=3}}

\begin{remark}

Theorem~\ref{thm:forbidden-subgraph} is specific to $d=2$, and for the same
reason Chapter~\ref{ch:higher-d-witnesses}'s collapse was specific to
$d\ge3$: the proof's forward direction rests on
Proposition~\ref{prop:tree-no-constraint} making trees exactly safe, which
in turn requires that individually-admissible edge data always extends —
true at $d=2$ only because $\mathrm{US}(2)=\mathrm{OS}(2)$
(Theorem~\ref{thm:n2degeneracy}) makes edgewise admissibility the same
condition over $\R$ and $\C$. At $d\ge3$, Chapter
\ref{ch:higher-d-witnesses} showed a single edge (a tree with one edge,
$\beta_1=0$) can already carry data admissible over $\C$ but not $\R$
(Theorem~\ref{thm:n3incompleteness}'s DFT example). So at $d\ge3$, the
topological condition "$G$ is a forest" is \emph{not} sufficient for
$\mathcal R_\R(G,d)=\mathcal D(G,d)$ — trees can already fail, so the
$d=2$ forbidden-subgraph theorem does not generalize by simply keeping the
same statement. What replaces it at $d\ge3$ is not a purely topological
condition at all, but the data-dependent single-edge obstruction of
Chapter~\ref{ch:higher-d-witnesses}: the relevant "forbidden pattern"
becomes a specific matrix (or class of matrices, per
Chapter~\ref{ch:unistochastic}'s open question about whether every
non-orthostochastic unistochastic matrix fails for the same reason), not a
graph shape.
\end{remark}

\begin{ledgerentry}[End of Chapter~\ref{ch:forbidden-subgraph}]
\textbf{Established:} the forbidden-subgraph question of
Chapter~\ref{ch:minimal-witnesses} is fully resolved at $d=2$
(Theorem~\ref{thm:forbidden-subgraph}): $G$ is always real-realizable if
and only if it is a forest, via a subdivision lemma proved algebraically
and cross-checked computationally at two cycle lengths.\\
\textbf{Not established:} any analogous purely topological
characterization at $d\ge3$, where the obstruction is data-dependent even
on trees (Chapter~\ref{ch:higher-d-witnesses}); the density-threshold
question of Chapter~\ref{ch:minimal-witnesses} remains entirely untouched
and would require a probability measure on edge-label data this book still
has not defined.
\end{ledgerentry}
